\scanpage{80}
\mspage{29}

La stabilité par chgt. de base est immédiate sur la forme c). \deleted{\unsure{La}} Stabilité par composition moins triviale : \diagram{\begin{tikzcd}[ampersand replacement=\&,row sep=small,column sep=small] {x\quad X}\arrow[d,"f"]\&{{}_{x}\backslash X}\arrow[l,Rightarrow]\arrow[d]\&\overline X\arrow[l]\arrow[d]\\ {y\quad Y}\arrow[d,"g"]\&{{}_{y}\backslash Y}\arrow[l]\arrow[d]\&{\deleted{\unsure{Z}}\overline Y}\arrow[l]\arrow[d]\\ {z\quad Z}\&{{}_{z}\backslash Z}\arrow[l]\&\{u\}\arrow[l] \end{tikzcd}} Dans le diagramme à carrés cart. ci-contre (déduit de $x\in\mathrm{Ob}\,X$, d'où $y=f(x)$, $z=g(y)=gf(x)$\deleted{\unsure{)}}, et de $u\in\mathrm{Ob}\,{}_{z}\backslash Z$), il faut prouver que $\overline X$ est $W$-asphérique. On sait que $\overline Y$ l'est ($g$ propre) et que $\overline X\mathrel{\relbar}\overline Y$ est à fibres $W$-asphériques. Il suffit donc d'utiliser la prop. 5 ci-dessous, et la remarque immédiate :

\uline{Lemme} Si $f$ est propre, il en est de $=$ de \deleted{qui est} $ {}_{x}\backslash f:{}_{x}\backslash X\longrightarrow {}_{y}\backslash Y$, \inserted{pour tt $x$ de $X$, $y=f(x)$,} En effet, si $u'$ est dans $ {}_{x}\backslash X$ et $v'$ dans $ {}_{y}\backslash Y$, \deleted{alors} et si $x',y'$ sont leurs images sur $X$, $Y$ resp., alors \[{}_{u'}\backslash({}_{x}\backslash X)\simeq {}_{x'}\backslash X,\qquad {}_{v'}\backslash({}_{y}\backslash Y)\simeq {}_{y'}\backslash Y\quad\deleted{\text{et}}\] et la flèche entre les deux s'identifie à la flèche de localisation \[{}_{x'}\backslash X\longrightarrow {}_{y'}\backslash Y,\]

% slanted
\margin{\deleted{En fait, si $f$ propre, $f$ est asphérique, \unsure{si} $f$ lisse, $f$ est \unsure{coasph.}}} \uline{Prop 5} Soit $f:X\longrightarrow Y$ de $\mathrm{Cat}$, supposons $f$ $W$-propre \deleted{et} $W$-lisse, et à fibres $W$-asphériques, alors $f\in W$ (et $=$ $f$ est \uline{universellt} dans $W$, i.e. il est dans $W$ et le reste par \uline{tout} chgt. de base).

% archive sheet number: 80
